\subsection{Prompting, Decoding, and Answer Scoring}
\label{app:scoring_protocol}

All evaluations use native chat templates, BF16 computation, eager attention, and one-beam greedy decoding, with Qwen's optional thinking mode disabled. Each prompt contains a single user message without demonstrations. Contextual QA prompts request the answer alone in its original spelling. Question-token positions are identified from the question stem's character offsets in the native template. Batched inputs are left-padded, and responses are trimmed at the first native end-of-sequence (EOS) token.

For \dataset\ and HotpotQA, the generation budget is 32 tokens. Primary accuracy requires both EOS completion and normalized exact match against an accepted answer alias. Following the 2Wiki-MultihopQA/HotpotQA evaluator, normalization lowercases text, removes punctuation and articles, and collapses whitespace \citep{ho2020constructing,yang2018hotpotqa,alab2020twikievaluation}. We also report text exact match without the EOS condition, token F1, and completion rates (\cref{tab:scoring}). Completed-answer and text exact match coincide for the reported composed-question results.

The composition gap persists under permissive alias matching. Among the 234/236/243 failed compositions with two correct constituent answers, 14/9/6 responses contain an accepted normalized alias as a complete word sequence. Crediting these responses yields failure rates of 48.1/49.1/47.3\% on the same 457/462/501-source subsets.

For MMLU, we select the highest-probability native token among A/B/C/D at the prefill used to generate the first response token \citep{hendrycks2021mmlu}. Each label is verified to be a distinct single token. Primary accuracy is micro-averaged over all 14,042 questions; generated-letter and subject-macro accuracy provide secondary diagnostics. GSM8K prompts request step-by-step reasoning followed by a final nonempty line of the form \texttt{\#\#\#\# NUMBER} \citep{cobbe2021gsm8k}, with a generation budget of 512 tokens. After removing thousands separators, correctness is determined by decimal numerical equality. Answer validity is assessed from the final line independently of EOS, with format failures and length cutoffs recorded separately. Every method uses the same zero-shot protocol for each benchmark.

\paragraph{Text accuracy and answer completion.} \method\ increases composed text F1 from 39.48/41.87/49.83\% to 73.41/74.35/76.85\%, gains of 33.93/32.48/27.02 points. Every \method\ composed response completes within the 32-token budget. \Cref{tab:scoring} reports exact match, partial-match F1, and completion for Native, full fine-tuning, and \method.

\begin{table}[H]
  \centering
  \caption{Scoring on 1,000 composed test questions (percent; adapted methods averaged over runs). Text EM removes the EOS condition; F1 uses the same aliases.}
  \label{tab:scoring}
  \small
  \setlength{\tabcolsep}{5pt}
  % Generated by analysis/build_tables.py.
\begingroup\renewcommand{\arraystretch}{1.08}
\begin{tabular}{llrrrr}
\toprule
\textbf{Model} & \textbf{Method} & \textbf{Completed EM} & \textbf{Text EM} & \textbf{Text F1} & \textbf{EOS completion} \\
\midrule
\rowcolor{tablebase}Qwen3-4B & Native & $26.90$ & $26.90$ & $39.48$ & $96.50$ \\
 & Full fine-tuning & $57.13$ & $57.13$ & $65.26$ & $99.93$ \\
 & Evidence filtering & $47.23$ & $47.23$ & $58.66$ & $97.43$ \\
 & Head tuning & $60.13$ & $60.13$ & $67.66$ & $\mathbf{100.00}$ \\
 & Scoped adaptation & $60.43$ & $60.43$ & $67.76$ & $\mathbf{100.00}$ \\
 & Graph-first selection & $63.37$ & $63.37$ & $69.92$ & $\mathbf{100.00}$ \\
 & Attention-only reader & $\mathbf{66.80}$ & $\mathbf{66.80}$ & $\mathbf{73.43}$ & $\mathbf{100.00}$ \\
 & Metadata-only reader & $64.03$ & $64.03$ & $70.83$ & $\mathbf{100.00}$ \\
\rowcolor{tableours} & \method\ (Ours) & $\mathbf{66.80}$ & $\mathbf{66.80}$ & $73.41$ & $\mathbf{100.00}$ \\
\midrule
\rowcolor{tablebase}Qwen3-8B & Native & $29.20$ & $29.20$ & $41.87$ & $98.30$ \\
 & Full fine-tuning & $55.17$ & $55.17$ & $63.55$ & $99.93$ \\
 & Evidence filtering & $48.10$ & $48.10$ & $59.66$ & $98.10$ \\
 & Head tuning & $61.80$ & $61.80$ & $69.44$ & $\mathbf{100.00}$ \\
 & Scoped adaptation & $60.90$ & $60.90$ & $68.61$ & $\mathbf{100.00}$ \\
 & Graph-first selection & $63.70$ & $63.70$ & $70.92$ & $\mathbf{100.00}$ \\
 & Attention-only reader & $\mathbf{67.20}$ & $\mathbf{67.20}$ & $\mathbf{74.35}$ & $\mathbf{100.00}$ \\
 & Metadata-only reader & $64.60$ & $64.60$ & $71.75$ & $\mathbf{100.00}$ \\
\rowcolor{tableours} & \method\ (Ours) & $\mathbf{67.20}$ & $\mathbf{67.20}$ & $\mathbf{74.35}$ & $\mathbf{100.00}$ \\
\midrule
\rowcolor{tablebase}Llama-3.1-8B & Native & $36.40$ & $36.40$ & $49.83$ & $\mathbf{100.00}$ \\
 & Full fine-tuning & $35.20$ & $35.20$ & $44.29$ & $98.73$ \\
 & Evidence filtering & $53.43$ & $53.43$ & $64.97$ & $\mathbf{100.00}$ \\
 & Head tuning & $63.60$ & $63.60$ & $70.70$ & $\mathbf{100.00}$ \\
 & Scoped adaptation & $66.30$ & $66.30$ & $73.59$ & $\mathbf{100.00}$ \\
 & Graph-first selection & $65.73$ & $65.73$ & $72.96$ & $\mathbf{100.00}$ \\
 & Attention-only reader & $\mathbf{69.70}$ & $\mathbf{69.70}$ & $76.78$ & $\mathbf{100.00}$ \\
 & Metadata-only reader & $66.70$ & $66.70$ & $74.03$ & $\mathbf{100.00}$ \\
\rowcolor{tableours} & \method\ (Ours) & $\mathbf{69.70}$ & $\mathbf{69.70}$ & $\mathbf{76.85}$ & $\mathbf{100.00}$ \\
\bottomrule
\end{tabular}
\endgroup

\end{table}

\subsection{Selected-Head Adaptation and Fixed Restoration}
\label{app:training_protocol}

Model adaptation uses the 2,000 composed training questions. The teacher-forced objective averages cross-entropy over the gold answer and native EOS, excluding prompt tokens. We optimize zero-initialized FP32 deltas for the full-rank Q/K/V/O slices of the selected heads while freezing the base parameters. One pass through the fixed source order uses single-example microbatches, eight-example gradient accumulation, and 250 updates. AdamW \citep{loshchilov2019adamw} uses a constant learning rate of $2\times10^{-5}$, $(\beta_1,\beta_2)=(0.9,0.999)$, $\epsilon=10^{-8}$, zero weight decay, and global gradient-norm clipping at 1. The model remains in evaluation mode during adaptation, with activation checkpointing and no dropout.

\Cref{tab:parameters} reports the selected parameter counts. Projection restoration reinstates native V/O contributions at the prefill positions and layers specified in \cref{tab:scopes}, while retaining the other learned contributions. These model-specific settings were fixed before the reported evaluation suite. All component interventions share the same trained checkpoint.

\begin{table}[H]
  \centering
  \caption{Selected-head and full-reader sizes. Trainable weights count Q/K/V/O slice deltas; model fractions exclude the separately fitted reader.}
  \label{tab:parameters}
  \small
  \setlength{\tabcolsep}{5pt}
  % Generated by analysis/build_tables.py.
\begin{tabular}{lrrrr}
\toprule
Model & Selected heads & Reader weights & Trainable weights & Model fraction \\
\midrule
Qwen3-4B & 20 & 42 & 23,592,960 & 0.587\% \\
Qwen3-8B & 20 & 42 & 39,845,888 & 0.486\% \\
Llama-3.1-8B & 40 & 82 & 75,497,472 & 0.940\% \\
\bottomrule
\end{tabular}

\end{table}

\begin{table}[H]
  \centering
  \caption{Fixed native-projection restoration rules, with zero-based layer indices. Restoration affects selected-head contributions at the specified prefill positions.}
  \label{tab:scopes}
  \small
  \setlength{\tabcolsep}{6pt}
  \begin{tabular}{llll}
    \toprule
    \textbf{Model} & \textbf{Projections restored} & \textbf{Layers} & \textbf{Token positions} \\
    \midrule
    Qwen3-4B & V and O & 16--19 & Prefix before the question \\
    Qwen3-8B & V and O & 12--19 & Question prefill \\
    Llama-3.1-8B & V & All selected layers & Prefix before the question \\
    \bottomrule
  \end{tabular}
\end{table}

\subsection{Fitting and Applying the Paragraph-Address Readout}
\label{app:readout_protocol}

Each selected head contributes two paragraph features: the logarithm of the mean question-token attention mass and the logarithm of the mass at the final question token. Masses are floored at $10^{-8}$ before taking logarithms. Two metadata features encode the log paragraph-to-prompt token ratio and whether the paragraph title appears in the question. The title feature uses normalized substring matching for Qwen and normalized whole-word matching for Llama; candidate graphs use whole-word matching for both families. Full readers contain 42/42/82 weights. Attention-only readers use the head features, while metadata-only readers use the two metadata features.

Standardization and fitting use the 2,000 composed training examples. Features are standardized across all training paragraphs, with population standard deviations floored at $10^{-4}$. The zero-initialized linear scorer is optimized for 500 full-batch Adam steps \citep{kingma2015adam}, using learning rate 0.03 and penalty $0.01\lVert w\rVert_2^2$. The cross-entropy objective identifies the annotated bridge among the original paragraphs. We omit a shared intercept because it cancels in the paragraph-wise softmax. Native and adapted models each receive a separately fitted, state-aligned reader. At inference, the frozen scorer resolves ambiguous candidates derived from the input. Unique candidates require no attention prefill, and inapplicable graphs retain the original context.

\subsection{Full-Parameter Comparison and Training Cost}
\label{app:fullft_protocol}

Full fine-tuning updates all model parameters in training mode, using the same objective, source count, gradient accumulation, learning rate, and optimizer hyperparameters. It uses BF16 weights, CPU-offloaded FP32 optimizer state, and seed-dependent source permutations. This recipe is shared across all three models, and both training methods perform 250 updates. Under the recorded shared-resource schedule, head adaptation averages 552/1,566/747 seconds, compared with 4,007/34,916/41,512 seconds for full fine-tuning (\cref{tab:training_cost}). Answer-format and completion rates are reported in \cref{tab:secondary_qa}.

\begin{table}[H]
  \centering
  \caption{Training workload and measured duration across three recorded runs. Both methods use one pass over 2,000 composed sources.}
  \label{tab:training_cost}
  \small
  \setlength{\tabcolsep}{5pt}
  % Generated by analysis/build_tables.py.
\begingroup\renewcommand{\arraystretch}{1.08}
\begin{tabular}{llrrrr}
\toprule
\textbf{Model} & \textbf{Training} & \textbf{Sources} & \textbf{Updates} & \textbf{Weights (M)} & \textbf{Seconds, mean (range)} \\
\midrule
Qwen3-4B & Full fine-tuning & 2000 & 250 & 4022.47 & 4007 (3970--4031) \\
\rowcolor{tableours} & \method\ adaptation (Ours) & 2000 & 250 & 23.59 & 552 (496--656) \\
\midrule
Qwen3-8B & Full fine-tuning & 2000 & 250 & 8190.74 & 34916 (32817--37643) \\
\rowcolor{tableours} & \method\ adaptation (Ours) & 2000 & 250 & 39.85 & 1566 (1361--1787) \\
\midrule
Llama-3.1-8B & Full fine-tuning & 2000 & 250 & 8030.26 & 41512 (29470--47703) \\
\rowcolor{tableours} & \method\ adaptation (Ours) & 2000 & 250 & 75.50 & 747 (722--768) \\
\bottomrule
\end{tabular}
\endgroup

\end{table}

\subsection{Recorded Runs and Statistical Uncertainty}
\label{app:uncertainty}

The three recorded runs use seed identifiers 42/43/44. With fixed source order, zero initialization, and evaluation mode, selected-head adaptation produces identical loss and gradient-norm sequences across runs; reader fitting is also deterministic. Reported sample standard deviations summarize variation across the recorded evaluations, whose individual composed scores appear in \cref{tab:per_run}.

For native and selected-head contextual QA evaluations, Qwen3-4B and Llama use batch size 1 in run 42 and 4 in runs 43/44; Qwen3-8B uses 4 throughout. MMLU also uses batch size 4 throughout. GSM8K uses 1/1/4 for Qwen3-4B, 4/4/4 for Qwen3-8B, and 1/1/1 for Llama. Full-parameter evaluations all use batch size 4. BF16 greedy decoding can produce small batch-dependent differences. Qwen3-8B's main native composed accuracy is 29.2\%, compared with 28.9\% in its batch-1 access and necessity experiments. Each causal contrast is paired with the native reference from the same experiment.

Confidence intervals in \cref{tab:paired_effects} use paired percentile bootstraps with 2,000 replicates and random seed 42. For repair comparisons, we average each source's correctness over the three runs and then resample source groups defined by transitively shared normalized root/bridge titles. Each of the 1,000 test groups contains one source. The resulting intervals quantify source-sampling variation on the fixed panel. Factorial contrasts instead resample the 25 shared-context quadruplets, preserving all four receivers and their paired intervention arms within each group.

\begin{table}[H]
  \centering
  \caption{Paired composed-accuracy changes in percentage points, with 95\% bootstrap intervals. Repair comparisons use source-wise run averages; bridge calibration uses its own experiment's native reference.}
  \label{tab:paired_effects}
  \small
  \setlength{\tabcolsep}{5pt}
  % Generated by analysis/build_tables.py.
\begingroup\renewcommand{\arraystretch}{1.08}
\begin{tabular}{lrrr}
\toprule
\textbf{Paired contrast} & \textbf{Qwen3-4B} & \textbf{Qwen3-8B} & \textbf{Llama-3.1-8B} \\
\midrule
\method\ $-$ Native & $39.90\;[36.73, 43.03]$ & $38.00\;[34.70, 41.40]$ & $33.30\;[30.23, 36.50]$ \\
Full readout $-$ paragraph order & $3.43\;[2.40, 4.63]$ & $3.50\;[2.40, 4.80]$ & $3.97\;[2.87, 5.30]$ \\
Bridge calibration $-$ Native & $4.90\;[2.90, 6.80]$ & $5.50\;[3.70, 7.40]$ & $16.00\;[13.40, 18.60]$ \\
\bottomrule
\end{tabular}
\endgroup

\end{table}

\begin{table}[H]
  \centering
  \caption{Composed accuracy (percent) on the same 1,000 test sources for each recorded run. Interpretation of run dispersion follows Appendix~\ref{app:uncertainty}.}
  \label{tab:per_run}
  \small
  \setlength{\tabcolsep}{6pt}
  % Generated by analysis/build_tables.py.
\begingroup\renewcommand{\arraystretch}{1.08}
\begin{tabular}{llrrr}
\toprule
\textbf{Model} & \textbf{Method} & \textbf{Run 42} & \textbf{Run 43} & \textbf{Run 44} \\
\midrule
Qwen3-4B & Full fine-tuning & 54.4 & 58.6 & 58.4 \\
\rowcolor{tableours} & \method\ (Ours) & 66.4 & 67.0 & 67.0 \\
\midrule
Qwen3-8B & Full fine-tuning & 52.7 & 61.3 & 51.5 \\
\rowcolor{tableours} & \method\ (Ours) & 67.2 & 67.2 & 67.2 \\
\midrule
Llama-3.1-8B & Full fine-tuning & 32.8 & 33.7 & 39.1 \\
\rowcolor{tableours} & \method\ (Ours) & 69.5 & 69.8 & 69.8 \\
\bottomrule
\end{tabular}
\endgroup

\end{table}
